\section{Certificates, research coverage, and the remaining frontier}\label{sec:certification}
\subsection{Automated allocation and verification}\label{sec:verification}
The computation has three separate layers: the frozen source census,
a discovery pipeline, and an exact categorification-certificate checker. Discovery is permitted to time
out or to use numerical linear programming to propose a functional.
Only the checker can validate a negative record. The final ledger
contains one of four disjoint reporting classes: generally excluded,
unitarily realized, spherically realized with no pseudo-unitary
categorification, or unresolved. Five separate existence columns record
the categorical structures, with no implication from spherical to
pseudo-unitary existence. The pipeline rejects any intersection
of its positive and negative sets.

\begin{table}[p]\centering\small
\setlength{\tabcolsep}{3pt}
\begin{tabular*}{\textwidth}{@{\extracolsep{\fill}}lrrrrrr@{}}\toprule
Criterion & $m=1$ & $m=2$ & $m=3$ & $m=4$ & Total removed & Seconds\\\midrule
Input & 39 & 154 & 384 & 872 & --- & ---\\
Rank-two subring & 0/39 & 4/150 & 4/380 & 49/823 & 631 & 0.9\\
$d$-number & 10/29 & 87/63 & 312/68 & 687/136 & 14753 & 11.6\\
Zero spectrum & 2/27 & 0/63 & 1/67 & 0/136 & 4 & 0.6\\
Lower-rank subring & 0/27 & 0/63 & 6/61 & 16/120 & 144 & 0.4\\
Involution module & 4/23 & 3/60 & 19/42 & 13/107 & 233 & 52.0\\
Dimension quotient & 0/23 & 1/59 & 2/40 & 2/105 & 20 & 6.5\\
Cyclotomicity & 0/23 & 0/59 & 2/38 & 1/104 & 14 & 2.2\\
Morita descent & 0/23 & 2/57 & 1/37 & 1/103 & 4 & 0.6\\
Field automorphisms & 0/23 & 1/56 & 0/37 & 0/103 & 2 & 5.8\\
Weighted induction & 0/23 & 1/55 & 1/36 & 1/102 & 7 & 5.5\\
Sphericalization/induction & 2/21 & 33/22 & 32/4 & 85/17 & 693 & 857.6\\
Dimension order & 0/21 & 0/22 & 0/4 & 0/17 & 1 & 184.1\\
Character induction & 0/21 & 0/22 & 0/4 & 0/17 & 4 & 98.7\\
Tensor rotation & 0/21 & 0/22 & 0/4 & 0/17 & 0 & 81.8\\
Modular induction & 0/21 & 0/22 & 0/4 & 0/17 & 0 & 65.2\\
Deferred exact checks & 0/21 & 0/22 & 0/4 & 0/17 & 1 & 157.2\\
\midrule
Constructed models & 21 & 22 & 3 & 2 & --- & ---\\
Unresolved & 0 & 0 & 1 & 15 & --- & ---\\
\bottomrule\end{tabular*}
\par\vspace{1.3em}
\begin{tabular*}{\textwidth}{@{\extracolsep{\fill}}lrrrr@{}}\toprule
Criterion & $m=5$ & $m=6$ & $m=7$ & $m=8$\\\midrule
Input & 1582 & 2768 & 3814 & 6999\\
Rank-two subring & 81/1501 & 121/2647 & 113/3701 & 259/6740\\
$d$-number & 1356/145 & 2442/205 & 3473/228 & 6386/354\\
Zero spectrum & 0/145 & 1/204 & 0/228 & 0/354\\
Lower-rank subring & 25/120 & 28/176 & 22/206 & 47/307\\
Involution module & 49/71 & 23/153 & 78/128 & 44/263\\
Dimension quotient & 6/65 & 2/151 & 4/124 & 3/260\\
Cyclotomicity & 1/64 & 3/148 & 0/124 & 7/253\\
Morita descent & 0/64 & 0/148 & 0/124 & 0/253\\
Field automorphisms & 0/64 & 1/147 & 0/124 & 0/253\\
Weighted induction & 1/63 & 1/146 & 1/123 & 1/252\\
Sphericalization/induction & 58/5 & 140/6 & 115/8 & 228/24\\
Dimension order & 0/5 & 0/6 & 0/8 & 1/23\\
Character induction & 0/5 & 0/6 & 0/8 & 4/19\\
Tensor rotation & 0/5 & 0/6 & 0/8 & 0/19\\
Modular induction & 0/5 & 0/6 & 0/8 & 0/19\\
Deferred exact checks & 0/5 & 0/6 & 1/7 & 0/19\\
\midrule
Constructed models & 1 & 3 & 2 & 0\\
Unresolved & 4 & 3 & 5 & 19\\
\bottomrule\end{tabular*}
\caption{The executed ordered sieve. Criterion cells are newly excluded / remaining. Seconds are recorded pass wall times; 4 or 8 workers per stage are recorded in the ledger. Sphericalization/induction checks retained proof trees for $m\le7$ and discovers proofs at $m=8$. Trees can use nested advanced filters; unfinished retained checks appear in the deferred row. Older residuals are retained. Later dimension-order, character, tensor-rotation and modular discovery targets $m\in\{3,8\}$. Timeouts remove nothing. This practical order is not claimed optimal.}\label{tab:stages}
\end{table}


Tables~\ref{tab:stages} and~\ref{tab:stages-high} give a sequential view: a tensor removed at one
stage is not counted as a new removal at any later stage. The order
places inexpensive structural and arithmetic obstructions before
larger finite branch searches. Effectiveness means the number of
\emph{additional} tensors decided at that point, rather than the number
that a criterion could reject when run in isolation. The table is a
transparent operational order, not a comparison of logical strength or
a hardware-independent timing ranking. Ineligible hypotheses, an
unfinished search and a feasible necessary relaxation all leave a
ring undecided unless another accepted certificate or model resolves it.

\begin{table}[ht]\centering
\begin{tabular}{lr}\toprule
Retained proof family & Rings\\\midrule
Rank-two fusion subring & 631\\
Formal-codegree $d$-number obstruction & 14753\\
Zero-spectrum obstruction & 4\\
Complete lower-rank subring descent & 144\\
Involution-module representation obstruction & 235\\
Subcategory dimension divisibility & 20\\
Cyclotomicity & 14\\
Solvability descent and involution-return indicators & 4\\
Noncommuting dimension-field automorphisms & 2\\
Multiplicity-weighted noncommutative induction & 7\\
Commutative signed sphericalization & 192\\
Complete oriented signed sphericalization & 376\\
Direct general-to-positive reduction & 70\\
Further general-to-positive certificates & 57\\
Signed reduction with the common induction test & 2\\
\midrule
Total & 16511\\
\bottomrule\end{tabular}
\caption{Disjoint retained proof families, listed in the order of the corresponding sieve stages. Several signed and positive-dimension families occur within the sphericalization/induction stage. Nested criteria are not counted a second time.}\label{tab:allocation}
\end{table}


The allocation in Table~\ref{tab:allocation} is disjoint. It records
the certificate family stored for each exclusion, while
Tables~\ref{tab:stages} and~\ref{tab:stages-high} make the sequential accounting explicit.
Neither table asserts mathematical independence of the criteria. The signed and positive-dimension records call the
same low-induction and finite-indicator framework internally. When an
older unbalanced relative-trace certificate suffices, it is interpreted
as a certificate for Theorem~\ref{thm:balanced} with zero coefficients
on the added balance equations. Likewise every inherited Farkas vector remains valid after the Galois-forbidden columns of Theorem~\ref{thm:character-induction} are removed: the inequalities on retained columns are unchanged. The stronger framework therefore retains the old exact witnesses rather than treating their special cases as unrelated obstructions.

Every certificate is verified from the original ordered tensor. The
checker reconstructs number fields, verifies separating elements and
factorizations, recomputes low-column and signed-table possibilities,
and checks all integer inequalities. In particular it does not assume
that a stored list of branches is exhaustive. For Farkas records it
checks \eqref{eq:farkas} with Python integers. During certificate
verification the floating-point linear-programming entry point is
replaced by a function that raises an error if called.

The source revision and tensor hashes fix the counting input. No census
is regenerated for the categorification analysis. The positive check
reconstructs the listed products, the zesting tensor and scalar
coboundary, finite-group character rings, and group-theoretical
convolution tensors, including the twisted $S_4$ cochains and their
projective stabilizer factors. Exact Weyl characters reconstruct the
$G_2$ model. Every positive entry, including every multiplicity-one
entry, is matched by an explicit unit-fixing basis permutation; the
complete $\ConstructorCount$-record constructor catalogue is checked for omissions. Regression tests compare the native integer
routines with straightforward exact implementations, and deliberately
corrupted certificates must be rejected.

This is an exact computer-assisted verification, not a formalization in
a proof assistant. The categorical input theorems, their hypotheses,
and the mathematical completeness argument remain part of the written
proof. A checksum proves file integrity, not a mathematical theorem.

\begin{proof}[Proof of Theorem~\ref{thm:classification}]
Each negative record is an application of Theorems~\ref{thm:involution}
and~\ref{thm:automorphisms}, or of Proposition~\ref{prop:inputs},
Corollary~\ref{cor:subring}, or of an exhaustive signed/positive reduction
followed by Theorem~\ref{thm:weighted} or Theorem~\ref{thm:balanced}
and the cited rank-four or rank-five tests, with the tensor-sign
reduction of Lemma~\ref{lem:positive-sign} and the extension of
Theorem~\ref{thm:rotation} where recorded. The stronger character-aware system of Theorem~\ref{thm:character-induction}, dimension-order test of Proposition~\ref{prop:dimension-order}, and signed refinement of Theorem~\ref{thm:signed-traces} are used where recorded. The spectral refinement of Proposition~\ref{prop:module-perron},
rank-three subring descent, Theorem~\ref{thm:morita-order6}, and
Proposition~\ref{prop:biquadratic-trace-obstruction} supply the
additional continuation certificates. The second continuation also
uses the ramification criterion of Proposition~\ref{prop:ramification-trace},
Theorem~\ref{thm:formal-codegree-trace-ideal}, the signed saturation
corollary, and Proposition~\ref{prop:quadratic-gauss-family}. The higher
stabilizer and ambient cyclotomic refinements use
Corollary~\ref{cor:stabilizer-average} and
Theorem~\ref{thm:quartic-stabilizer}, with a complete general-to-pseudo-unitary
reduction whenever their inner certificate has only pseudo-unitary scope.
Theorem~\ref{thm:stabilizer-rotation-frontier} supplies the final general
multiplicity-three exclusion: its sphericalization reduction, full
Morita-action calculation, local rotation lemma and complete exact
indicator separation are all required. The exact checker validates
all hypotheses and exhausts the required finite alternatives. Thus no
excluded tensor has a complex categorification. Section~\ref{sec:positive}
constructs and tensor-matches all claimed positive examples, with the
additional models in Sections~\ref{sec:literature-models} and~\ref{sec:g2-e4}, the higher models of Section~\ref{sec:through7-models}, and complete
catalogue in Section~\ref{sec:model-catalogue}. In particular the spherical
existence and pseudo-unitary nonexistence for $\RR{2}{124}$ use the
full-tensor identification and the two cited results of Schopieray. No complement argument is used
to assign the positive records. The census and disjoint ledger establish the counts in
Table~\ref{tab:classification}. The other records are left undecided,
with no assertion of nonexistence or realizability.
\end{proof}

\subsection{Regression controls}\label{sec:frontier}
The positively realized ring $\RR{2}{4}$ supplies independent controls
for the balanced and rotation systems. Their exact witnesses are
\filename{data/frontier_feasible.json} (formal rank $22$) and
\filename{data/frontier_rotation_feasible.json} (formal rank $27$).
Each is checked only against its own reconstructed necessary system;
the different formal ranks do not describe a single center. Existence
comes from Proposition~\ref{prop:g2-e4}.

Conversely, the generally excluded ring $\RR{3}{215}$ has an exact
rank-$22$ induction, fusion and modular-data relaxation. The stored
replay programs in \filename{research/m8_frontier/} preserve this
negative control for overinterpreting feasible data. The additional
stabilizer-rotation obstruction, not a failure of those weaker systems,
proves its nonexistence. These controls are kept outside the main
classification argument.

\subsection{Reproduction and counting data}\label{sec:reproduction}
The top-level command
\begin{center}\texttt{python verify.py --jobs 2}\end{center}
checks the frozen input correspondence, positive models, and allocated
negative certificates. The command with \texttt{python -O} checks
that correctness does not depend on Python assertions. Regression
controls exercise exact arithmetic, corrupted certificates and the
frontier witnesses. Precise dependencies, commands, report scope and
search-stop behavior are given in the accompanying verification guide.
These are classification checks, not a new enumeration of fusion rings.

The \filename{oeis/} directory contains an eight-term b-file, a CSV,
and proposed definition and provenance text. The definition is the
number of based-isomorphism classes of rank-six fusion rings of
\emph{exact} multiplicity $m$, without assuming commutativity,
self-duality, integrality or categorifiability. The terms are copied
from the pinned census. No accession number is assumed and no online
submission is claimed. Counting farther and classifying farther are
separate achievements.

\subsection{Research coverage and the next unresolved stratum}\label{sec:research-coverage}
The residual list is derived from the complete decision ledger, not
from a hand-selected sample. Multiplicities one through three are completely decided.
The smallest unresolved multiplicity and the exact residual tensors
are reported in Section~\ref{sec:results}; the first such tensor is the
priority for further theoretical and automated work. The stage table
also shows whether difficulties reappear at larger multiplicities.
A criterion applied successfully at one bound remains the same
mathematical criterion at the next bound: additional computation alone
is not presented as a new theorem.

The module-Gram refinements and relative-center constructions are
necessary tests or model constructions with explicit hypotheses.
A failed eligibility check does not exclude a ring. An inconclusive
Gram factor does not supply a model. For the relative Fibonacci center
of $TT_3$, the categorical construction and exact tensor reconstruction
supply the positive conclusion; feasible Gram data by themselves do not.

Literature comparisons distinguish exact fusion-rule matches from
shared dimensions. Unitary-only classifications do not exclude arbitrary
complex fusion categories, and candidate modular data are not
categorifications unless the source supplies a construction. The same
principles govern new models and new necessary criteria: a model must
have a justified categorical source and an exact full-tensor match;
an obstruction must state its hypotheses and exhaust all branches
needed for the claimed categorical structure.

The certificate files and immutable campaign manifests preserve provenance
without entering the mathematical criteria as lists of forbidden IDs.
The authoritative ledger is rebuilt from their disjoint union. Alternative
proofs can be retained, but a tensor receives only one exclusion allocation.
\filename{software/continuation_controls.py} independently checks accepted
continuation records with numerical optimizers disabled; \filename{verify.py}
replays the full current ledger. Discovery logs and incomplete searches
are not verification receipts.

The present frontier is described by complete tensor-derived necessary
data in the current research report. Its counts are not predictions about
categorifiability beyond multiplicity eight. A surviving integral profile,
fusion tensor or modular datum is assigned categorical existence only when
a categorical construction and exact tensor identification are supplied.

\subsection{Exact completion of the quartic centralizer data}
\label{subsec:frontier-quartic-completion}

The full tensor of $R_{6,4,91}$ contains the Ising based subcategory
with $g=b_4$ and $\sigma=b_3$. The preceding signed and induction
reductions leave two possible center ranks, $36$ and $42$. Each profile
has exactly one central simple restricting to $g$ and exactly one
restricting to $\sigma$. Their products, determined by restriction and
simplicity, give an Ising fusion subcategory of the center. Every braided
Ising category is nondegenerate. Proposition~\ref{prop:central-packets-adjunction}
therefore applies to its centralizer $\mathcal E$.

\begin{theorem}[Complete centralizer fusion and modular data for $R_{6,4,91}$]
\label{thm:91-centralizer-completion}
If $R_{6,4,91}$ has a pseudo-unitary categorification $\mathcal C$, then
$Z(\mathcal C)\simeq\mathcal I\boxtimes\mathcal E$, where $\mathcal I$
is Ising and $\mathcal E$ has rank $14$. The center-rank-$36$ profile is
impossible. For the rank-$42$ profile, the complete compatible fusion
ring of $\mathcal E$ has one based-ring isomorphism class. With its
restriction columns labelled as in the certificate, it has eight fusion
tensors, all related by permutations of repeated restriction columns.
There are exactly two compatible pairs of centralizer modular data;
they are complex conjugate. These assertions specify necessary central
data and do not assert that either pair is categorifiable.
\end{theorem}
\begin{proof}
Write the low summands of canonical unit induction as $A_0,\ldots,A_5$.
Their known character eigenvalues at the pure $g$ and $\sigma$ columns
identify their Ising coordinates without choosing a braiding on
$\mathcal C$:
\[
 A_0=1\boxtimes1,\quad A_1=g\boxtimes a,\quad
 A_2=g\boxtimes b,\quad A_3=1\boxtimes c,\quad
 A_4=\sigma\boxtimes u,\quad A_5=\sigma\boxtimes u^*.
\]
Both inherited low character labellings give these same coordinates
and the same duality permutation $(4\ 5)$ on low labels. Thus the right
adjoint of $F|_{\mathcal E}$ at the tensor unit is $1\oplus c$.

The complete central profile must be a union of packets
$(v,gv,\sigma v)$. The nonnegative integral packet-cover equations
have one solution for each of the two profiles. Each admits $144$
endpoint orientations up to permutations of equal columns. Of these,
$135$ violate the displayed low Ising coordinates. For a remaining
orientation, let $V$ have the proposed centralizer restriction columns.
The adjunction identity fixes
\[
 N_c=V^{\mathsf T}V-I.
\]
Nonnegative entries, the unit column, and the tensor-functor equation
$VN_c=N_{F(c)}V$ reject eight further orientations. Exactly one
restriction multiset survives for each rank. This is a decomposition
of the center; no tensor-product decomposition of $\mathcal C$ is used.

We describe the entire finite fusion search. For every possible duality
involution consistent with the restriction columns and the forced pair
$u,u^*$, use the symmetric variables
\[
 x_{ijk}=N_{ij}^{k^*},\qquad x_{ijk}=x_{\pi(i,j,k)}.
\]
The unit and $c$ products are fixed. Restriction of every product gives
linear equations and finite upper bounds
\[
 N_{ij}^{k}\le
 \min_{h:(v_k)_h>0}
 \left\lfloor\frac{(v_iv_j)_h}{(v_k)_h}\right\rfloor.
\]
Duality and commutation with $N_c$ give further exact linear equations.
Project this system onto multiplication by the simple whose restriction
is $b_1$, eliminating the complementary variables by rational row
reduction. Every actual fusion matrix is normal, since multiplication
by its dual is its transpose and the centralizer ring is commutative.

The rank-$12$ centralizer has one possible duality. Its projected affine
system has eight free coordinates and no bounded integral solution;
complete enumeration visits $11$ nodes. For rank $14$, there are four
possible dualities. Three give $64$ projected matrices each, none normal.
The remaining duality gives $512$ projected matrices, of which $16$ are
normal. For each such matrix, its commutation with all the other fusion
matrices is now linear. Eight completions are inconsistent; each of the
other eight has one bounded integral completion. Direct multiplication
checks the unit, duality, nonnegative integral coefficients and every
associativity identity. Swapping the three pairs of identical restriction
columns identifies these eight tensors in one based-ring orbit.

For each orbit representative the exact character table is obtained by
factoring a separating multiplication matrix over $\QQ(\zeta_{16})$.
All character products and all roots are checked exactly. In any modular
labelling, the character attached to object $j$ has formal codegree
$\operatorname{FPdim}(\mathcal E)/d_j^2$, and
$\widehat S_{ij}=d_j\chi_j(i)$ must be symmetric. Exhaustive assignment
of the characters subject to these requirements gives exactly two
$\widehat S$ matrices.

Put $z=\theta_\sigma$. The inherited second indicator is
$\nu_2(\sigma)=-1$, so
$z=\zeta_{16}^{e}$ with $e\in\{3,5,11,13\}$.
The low summands have twist one, hence
\[
 \theta_a=\theta_b=-1,\qquad \theta_c=1,\qquad
 \theta_u=\theta_{u^*}=z^{-1}.
\]
For each $S$ and each of these four values of $z$, apply balancing
with one of these known-twist objects in the first coordinate:
\[
 \theta_i\theta_j\widehat S_{ij}
   =\sum_k N_{i^*,j}^{k}d_k\theta_k.
\]
These are linear equations for all fourteen twists. Every branch has
rank fourteen. Exactly two are consistent, with Ising exponents $5$
and $11$, respectively. The unique solutions have twists in
$\mu_{16}$ and are complex conjugate. In particular, no a priori bound
on the unknown twist orders is being substituted for an exhaustive
search: the linear equations determine the twists themselves.
Each solution passes every balancing equation,
$\widehat S^2=\operatorname{FPdim}(\mathcal E)C$, the full modular
relation, and the requirement that the product central charge of
$\mathcal I\boxtimes\mathcal E$ is one.

The record is
\filename{research/round5/quartic_pair/R6_4_91_center_reconstruction.json}.
The reconstruction and optimizer-free replay are implemented in
\filename{software/round5_quartic_pair.py}. They recompute every packet
cover, orientation, duality, affine projection, normality test, fusion
completion, based-ring orbit, character labelling and twist solution
from the frozen full-tensor frontier. Replay also checks the inherited
signed reduction and its exhaustive profiles when requested.
\end{proof}

The surviving centralizer has no nontrivial invertible object. Its formal
codegree polynomial is
\[
 (t-8)^2(t^2-128t+2048)^2(t^2-64t+512)^2
 (t^4-2816t^3+151552t^2-1572864t+4194304).
\]
For the representative with $z=\zeta_{16}^5$, in the certificate's
unit-first order, the twist exponents are
\[
 (0,10,12,8,14,14,6,6,4,2,8,0,11,11).
\]
The other representative negates these exponents modulo sixteen.
An associator, compatible braiding and the requisite induction algebra
are not supplied by fusion coefficients and modular matrices alone.
Thus the theorem eliminates one inherited profile and completely fixes
the remaining fusion and modular data, while leaving the categorical
realizability question explicit.

\subsection{Full regular-action completion for $R_{6,4,22}$}
\label{subsec:frontier-22-completion}

The remaining quartic ring $R_{6,4,22}$ is treated by the same
restriction--word Gram and regular-action criteria, without a centralizer
factorization. The inherited center ranks are $30$ and $36$.

\begin{proposition}[The rank-$30$ branch]
\label{prop:22-rank30}
The center-rank-$30$ induction profile of $R_{6,4,22}$ is impossible.
\end{proposition}
\begin{proof}
Apply the quadratic and cubic word Gram criterion and couple its coefficient
counts to the projected twist types by
Proposition~\ref{prop:projected-product-balancing}. For each of the two
inherited character labellings the simultaneous system has $2950$
equations and $4132$ nonnegative variables. Each has an exact rational
separating vector. The certificate
\filename{research/round5/fermion/R6_4_22_certificate.json} and
\filename{software/round5_22.py} reconstruct the full domains and check
both separating vectors over the rationals. No integrality or numerical
infeasibility assertion is needed for these two exclusions.
\end{proof}

\begin{theorem}[Complete necessary central data for $R_{6,4,22}$]
\label{thm:22-central-completion}
If $R_{6,4,22}$ has a pseudo-unitary categorification, its center has rank
$36$. Its complete central fusion ring, with the forgetful restriction
map, has one based-ring isomorphism class among all the inherited
possibilities. Each of the two conjugate low character labellings has
exactly one compatible modular datum. The two modular data are complex
conjugate and have twist order $20$.
\end{theorem}
\begin{proof}
For the rank-$36$ profile, retain all $50$ quadratic and cubic words in
the five nonunit low central simples. The complete conditional Gram
ellipsoids of Theorem~\ref{thm:joint-restriction-word-gram} give $155$
possible pairs consisting of a restriction column and its word
coefficient vector. Their simultaneous count system has $902$ equations.
Exact bounded affine enumeration has rank $88$, visits $136389$ nodes,
and gives precisely five nonnegative integral count vectors. The two
low labellings give identical word data, which is verified directly;
thus this count enumeration covers both.

Expand each count vector into the individual central simples. Duality
must dualize its restriction column and permute its word coordinates by
the known low duality. In each of the five patterns this leaves two
possible involutions, corresponding to the two indistinguishable copies
of one column. Apply Proposition~\ref{prop:word-action-completion} to
multiplication by $A_1$. In matrix row convention, its entries satisfy
\[
 (N_{A_1})_{jk}=(N_{A_1})_{k^*,j^*},\qquad
 U_wN_{A_1}=U_{wA_1},
\]
whenever both words have been retained. Tensoriality fixes their
restrictions, and the prescribed eigenvalues fix the full characteristic
polynomial. All variables have finite integral bounds from restriction
products. Exact interval propagation and affine enumeration reject both
dualities in four patterns, and one duality in the fifth pattern.

The last case has $1184$ affine matrices after the prescribed trace is
imposed. Exactly two have the required characteristic polynomial.
Adjoining the actions of $A_2,A_3,A_4$ successively, with commutation
against all previously fixed actions, gives one continuation of each;
$N_{A_5}=N_{A_4}^{\mathsf T}$. The resulting two complete low-action
systems differ only by swapping the indistinguishable central simples.
This equivalence is checked on every matrix entry.

There is a unique central simple whose restriction is $b_1$. Its
multiplication matrix is now determined by restriction, Frobenius
reciprocity and commutation with the fixed low actions: the exact
linear system has $465$ variables, rank $465$, and one nonnegative
integral solution. The central invertible and the other simple actions
are then completed in the same way. Each step has one solution.
The finished tensor is checked directly for the unit, duality,
commutativity, nonnegative integral coefficients and all associativity
identities. Hence the procedure has completed the entire central ring,
not only its products with low simples.

The exact character table of this ring is computed over
$\QQ(\zeta_{20})$ from a separating multiplication matrix. Formal
codegrees, symmetry and the prescribed low character eigenvalues leave
one $S$-labelling for each inherited labelling. Set the six low twists
equal to one. Balancing with these known-twist objects is a linear
system of rank $36$ in the $36$ unknown twists. It has a unique solution
for each labelling; the solutions are complex conjugate and belong to
$\mu_{20}$. The unique nonunit central invertible is forced to have twist $-1$.
Every balancing equation, the full
modular relation and central charge one are checked exactly.

The complete reconstruction is
\filename{software/round5_regular_completion.py}, using the generic
regular-action engine \filename{software/round5_low_actions.py}.
The record is
\filename{research/round5/quartic_pair/R6_4_22_center_reconstruction.json}.
Replay rebuilds the conditional Gram domains, all five count vectors,
all ten duality cases, every surviving matrix, the based-ring
identifications and both modular completions from the frozen full-tensor
frontier. It uses no numerical optimizer and remains valid with Python
assertions disabled.
\end{proof}

The two quartic rings have therefore reached complete necessary central
fusion and modular data. This is stronger than an integral induction or
Gram witness, but it still leaves categorical realization: a compatible
associator and braiding, together with the required canonical induction
algebra, have not been constructed. The current records retain this
remaining question explicitly.

\subsubsection{Finite checks after modular completion}

The completed matrices allow the same rotation criterion to be applied
with every central simple as a source, rather than only the central
unit. Write $S,T$ for the normalized, anomaly-one center matrices and
$W_{Zi}=[F(Z):X_i]$. The generalized indicator vector is
\[
 \bigl(\nu^{Z}_{m,1}(X_i)\bigr)_{Z,i}=ST^mW.
\]
The corresponding rotation on
$\operatorname{Hom}(F(Z),X_i^{\otimes n})$ has $n$th power
$\theta_Z^{-1}$. By \cite[Propositions~3.3--3.7]{BJT}, for
$g=\gcd(n,k)$ its normalized power trace is
\[
 \theta_Z^{k/n}\nu^Z_{n,k}(X_i)
 =\alpha_{k/g,n/g}\!\left(
     \theta_Z^{g/n}\nu^Z_{n/g,1}(X_i^{\otimes g})\right),
\]
where $\alpha_{a,m}(\zeta_m)=\zeta_m^a$. The expression to which this
Galois automorphism is applied must lie in $\QQ(\zeta_{n/g})$.
Exact Fourier inversion must give nonnegative integral eigenspace
dimensions whose sum is
$\dim\operatorname{Hom}(F(Z),X_i^{\otimes n})$.

For both conjugate modular data of each quartic ring, these field
membership and Fourier conditions hold for every original simple,
every central simple, and every $2\le n\le16$. The full center also
satisfies signed-permutation Galois covariance and the twist-square
compatibility. Finally,
\[
 \tau_n=\sum_Z d_Z^2\theta_Z^n>0
 \qquad (\gcd(n,\operatorname{ord}T)=1),
\]
for all eight coprime residue classes of each exponent. Thus every
coprime higher central charge is one, as required of a Drinfeld center
by \cite[Theorem~4.4]{NSWGauss}. Positivity is checked by exact real-root
isolation in the real quartic field. The equality
$\tau_n=\operatorname{FPdim}(\mathcal C)$ is not assumed for general
coprime $n$.

These additional surviving conditions are reconstructed in
\filename{software/round5_modular_checks.py}; its record is
\filename{research/round5/quartic_pair/modular_checks.json}.
They strengthen the verified necessary data and do not resolve the
remaining categorical realizability problem.


The saturated dimension-field characters of both quartic rings have order
four and are nonreal. Their conductors are respectively $20$ and $16$.
The real-character specialization of
Proposition~\ref{prop:saturated-modular-moments} is therefore not applied
to them. Instead, the exact audit
\filename{research/round5/quartic_pair/conductor_intertwiner.py}
constructs an equivariant quotient of
$\operatorname{Ind}_B^{\mathrm{SL}_2(\mathbb Z/f\mathbb Z)}\chi$
onto the cyclic span of $w_\chi$ in each completed modular datum.
It checks every $S,T$ generator edge and
$P_U S w_\chi=\lambda w_{\bar\chi}$, where $P_U$ projects onto
$T$-fixed vectors. The induced and cyclic dimensions are $(36,18)$
for $R_{6,4,22}$ and $(24,24)$ for $R_{6,4,91}$, for both saturated
characters and both modular variants. Complete reducibility and polar
decomposition give an isometric equivariant section of each quotient.
Thus every valid universal conductor moment or spectral identity already
holds on these cyclic spans, and adds no obstruction to the completed data.

The full central projection of
Proposition~\ref{prop:modular-subcategory-projection} also survives.
For the pointed rank-two subcategory of $R_{6,4,22}$ and the Ising
subcategory of $R_{6,4,91}$, twist and forgetful compatibility leave
respectively five and seven integral splitting variables in $B$.
The exact equations $S_{\mathcal D}B=BS_{\mathcal C}$ have ranks five
and seven and force one nonnegative integral matrix in each conjugate
modular variant. Every $S,T$, forgetful and unit equation is rechecked
in \filename{research/round5/quartic_pair/subcategory_projection.py}.
Thus the full projection condition adds no quartic obstruction.


\subsection{The generalized near-group ring $R^\prime(4)$}
\label{sec:frontier-155}
Write the basis of $\RR{4}{155}$ as $1,h,g,hg,X,Y$, where the
invertibles form $C_2^2$, $h$ fixes $X,Y$, and $g$ exchanges them.
The ring belongs to the family
\begin{equation}\label{eq:nonselfdual-neargroup}
 Y=X^*,\qquad XY=1+h+m(X+Y),\qquad X^2=g+hg+m(X+Y),
\end{equation}
with $m=4$. The following uniform argument eliminates signed-table
search for this family.

\begin{proposition}\label{prop:nonselfdual-neargroup-spherical}
Every complex fusion category with rules
\eqref{eq:nonselfdual-neargroup}, for $m\ge0$, is spherical and has a
pseudo-unitary Galois conjugate.
\end{proposition}
\begin{proof}
In the sphericalization choose the dimension-one lifts $k_+$ of the
invertibles, and choose $X_+$ and $Y_+=g_+X_+$. Their common nonzero
dimension $\delta$ forces
$h_+X_+=X_+$, $h_+Y_+=Y_+$ and $X_+^*=Y_+$.
The invertible part of the self-dual product $X_+Y_+$ is $1_++h_+$,
by reciprocity; hence
\[
 X_+Y_+=1_++h_++s(X_++Y_+)+t(X_-+Y_-),\qquad s+t=m.
\]
Thus $\delta^2=2+r\delta$ with $r=2(s-t)$; changing both lifts if
necessary gives $0\le r\le2m$. The squared-norm global dimension and
the Frobenius--Perron dimension are
\[
 d_r=8+r^2\pm r\sqrt{r^2+8},\qquad
 D_m=8+4m^2+2m\sqrt{4m^2+8}.
\]
The ratio $d_r/D_m$ is an algebraic integer
\cite[Proposition~8.22]{ENO}. For nonzero even $r$, $r^2+8$ is not a
square: $(v-r)(v+r)=8$ would force $r=1$. The same holds for $2m$.
When $m>0$, the compositum admits an automorphism negating both
quadratic roots (or only the denominator root when $r=0$).
Multiplying the ratio by its image gives the positive rational
algebraic integer
\[
 \frac{r^2+8}{4m^2+8}\le1.
\]
It equals one, so $r=2m$ and $t=0$. For $m=0$ this holds directly.
The lift $g_+$ commutes with $X_+,Y_+$, since both orders have the same
underlying simple and nonzero dimension. Therefore
$X_+^2=Y_+^2=g_+X_+Y_+$, and all products of the chosen lifts close
among them. They form a full spherical fusion subcategory projecting
tensor-equivalently onto the original category. Its noninvertible
dimension is $m\pm\sqrt{m^2+2}$; conjugating to the positive root
gives pseudo-unitarity.
\end{proof}

For $m=4$ the dimensions are
\[
 d=4+3\sqrt2,\qquad D=72+48\sqrt2.
\]
There is no nontrivial faithful grading: the adjoint subring is the
whole ring. The pointed subrings are its only proper nontrivial fusion
subrings. Thus neither a grading nor rank-three subring descent resolves
this case.

The retained central-induction certificate gives a finite necessary
domain of five profiles, with center ranks $46,48,52,54,60$.
For clarity, its categorical reductions are as follows. Pointed
subcategory projection for all three involutions leaves one low frame
and indicators $(\nu_2(X),\nu_2(Y),\nu_2(h),\nu_2(g),\nu_2(hg))
=(0,0,1,-1,-1)$. The $1+h$ module action and stabilizer-rotation
criterion constrain the third indicators. Exact Gram, Galois, word
and balancing equations then force central lifts of $h,g,hg$.
The lift of $h$ is fermionic: if it were bosonic, de-equivariantization
would split the fixed object $X$ into dimensions $d/2$, whose norm
is $-1/2$, contrary to integrality. A lift of $g$ generates a semion
subcategory of the center, which splits off from the center.

In its modular centralizer the restricted adjoint algebra is
$1+A_1+A_2$. For its restriction matrix $V$, adjunction gives
$V^{\mathsf T}V=I+N_{A_1}+N_{A_2}$. The two low simples have opposite
monodromy charges against the fermion. Consequently their multiplication
matrices are the same-charge and opposite-charge blocks of
$V^{\mathsf T}V-I$. Exhausting charge counts, duality, pointed actions,
restriction products and ribbon balancing leaves the six profiles,
each with a forced pair of low fusion matrices. For the rank-$44$ profile, the height-six full-character system of
Proposition~\ref{prop:round3-full-character}, in $\QQ(\zeta_{24})$,
has no solution with integral twist-type and rotation-eigenvalue counts.
The exact certificate gives $49$ integral-bound inferences followed
by one separating ray, using the bound argument in the proof of
Theorem~\ref{thm:1204-excluded}. Pointed restriction columns use their
proved twists $\pm1$ or $\pm i$; the other columns retain the complete
joint trace domain. All surviving character assignments differ only
by permutations of identical low columns, and no third-indicator value
is imposed. This removes the rank-$44$ profile.

The five remaining profiles are necessary conditions. The semion splitting of the center does
not assert a tensor-product splitting of the original category.

\begin{proposition}[No braided Morita representative]
\label{prop:155-no-braided-morita}
A categorification of $\RR{4}{155}$ is not Morita equivalent to any
braided fusion category.
\end{proposition}
\begin{proof}
Work with a pseudo-unitary Galois representative. Since the adjoint
subcategory is the whole category, the forgetful map on central
invertibles is injective. The central-induction certificate gives all four
central lifts, with twists $1,-1,i,-i$, up to complex conjugation.
Thus the center has no nonidentity invertible of twist one.

If a braided category $\mathcal M$ were Morita equivalent, it would
also be pseudo-unitary. The equivalence of centers preserves their
canonical positive ribbon structures. Every central invertible
forgetting to the unit of $\mathcal M$ has twist one, so this
forgetful map is injective. The image $a$ of the central semion has
its canonical lift supplied by the braiding of $\mathcal M$; uniqueness
forces that lift to be the semion itself. Hence $\langle a\rangle$
is nondegenerate, and braided decomposition gives
$\mathcal M\simeq\mathrm{Semion}\boxtimes\mathcal M_0$
\cite{DGNO}. Its center then contains the nonidentity twist-one
invertible $(a,a)$ from
$Z(\mathrm{Semion})\simeq\mathrm{Semion}\boxtimes\mathrm{Semion}^{\rm rev}$,
a contradiction. Galois conjugation preserves braidings and Morita
contexts, so the conclusion applies to the original category.
\end{proof}


\subsection{The multiplicity-four ring with a rank-35 central datum}
\label{subsec:frontier-629}

The following reduction determines the remaining modular and fusion data
for $R_{6,4,629}$. It does not construct its associators or a category.
All indices in the certificate files are zero based.

\begin{theorem}[Complete necessary modular reduction for $R_{6,4,629}$]
\label{thm:629-modular-reduction}
Every complex categorification of $R_{6,4,629}$ admits a spherical
structure and has a pseudo-unitary Galois conjugate. For that conjugate,
the center has rank $35$ and twist exponent $15$. Its restriction matrix,
normalized modular matrices and fusion tensor belong to one explicitly
computed class, up to complex conjugation and relabelling of simples
having identical restriction columns and twists.
The previously retained induction profiles of ranks $30$ and $38$ are
impossible. The rank-$35$ datum satisfies nonnegative integral Verlinde
multiplicities, but categorifiability remains unresolved.
\end{theorem}

\begin{proof}
The unrestricted signed-sphericalization enumeration has $32$ associative
quotients, all in the sign-gauge orbit of the original tensor. Changing
the chosen lifts by that gauge makes every opposite-lift multiplicity
zero, so the chosen lifts form a tensor subcategory equivalent to the
original category. Its spherical structure transports to the original
category. Of the two rational character components of the original
ring, the quadratic component has zero simple dimensions and cannot be
the spherical dimension character. The other component is one quartic
Galois orbit containing the Frobenius--Perron character. Conjugating to
that embedding gives pseudo-unitarity. The unrestricted enumeration and
the character calculation are exactly replayed by
\filename{research/round4/m4\_629/general\_reduction.json} and its
optimized replay.

The complete induction-profile reduction reconstructs all seven low
frames. Six fail the spectral condition. The remaining frame has two
compatible character labellings and $784$ integral high-column profiles.
Low-word Verlinde integrality and nonnegative restricted products leave
five profiles, of ranks $24,30,32,35,38$. The odd relative-Galois orbit
criterion and exact twist-moment separating certificates reject all five
profiles for one labelling and the ranks $24,32$ for the other. Thus the
input to the present computation is exhaustive: one labelling and three
profiles, with ranks $30,35,38$.

The six low restriction columns, in the retained order, are
\[
\begin{split}
 &(1,0,0,0,0,0),\quad (1,1,1,1,1,1),\quad
 (1,1,2,2,2,3),\\
 &(1,1,2,2,2,2),\quad (1,1,1,2,2,2),\quad
 (1,1,2,1,2,2).
\end{split}
\]
They are all self-dual, as follows from the real character table and its
low-character correspondence. No self-duality assumption is made about
the remaining central simples at this stage.

For each retained profile we rebuild the complete relative root-orbit
domain and the cyclic-rotation trace equations. If a column has
multiplicity $m$ and a relative orbit has degree $e$, its orbit-count
variable lies between $0$ and $\lfloor m/e\rfloor$. Every rotation
eigenmultiplicity lies between zero and the dimension of its rotation
space. These are finite integer intervals. Before rational row reduction
we put variables with larger upper bounds first. This is only a column
permutation, so it preserves the integer lattice and the solution set;
it makes the subsequent exact search much smaller. The surviving free
variables are enumerated with the inherited interval and integrality
pruning, and every leaf is substituted into the original equations.
The complete search results are
\[
\begin{array}{c|rrrrrr}
\text{central rank}&\text{height}&\text{equations}&\text{variables}&
 \text{linear rank}&\text{nodes}&\text{solutions}\\ \hline
30&6&4323&405&350&203&0\\
35&10&2476&587&523&4157&2\\
38&10&1942&527&470&2408&2
\end{array}
\]
Consequently the rank-$30$ profile is excluded. In each other profile the
two solutions are complex conjugate. Every occupied relative orbit has
degree one, so all twists are explicitly in $\mathbb Q(\zeta_{15})$.
The decoded twists have orders $3,5,15$, with least common multiple $15$.
This decoding is exact: a relative average of roots of unity having
absolute value one can only average identical roots. The decoded lists
are also checked to be conjugate as multisets of restriction--twist
pairs. The low simples are exactly the twist-one simples.

Let $D=\operatorname{FPdim}(\mathcal C)$ and write $d_Z$ for a central
simple dimension. The six normalized low rows are fixed by
\[
 S_{A_i,Z}=\lambda_i(Z)d_Z/D,
\]
where $\lambda_i(Z)$ is the joint-character eigenvalue attached to the
restriction column of $Z$. We apply
Proposition~\ref{prop:modular-word-closure}. Since the low objects are
self-dual and exhaust the twist-one objects, the low Galois signed
permutations are determined by unique matches in the known low block.
The congruence identities therefore supply the images of
$T^aSe_i$ for every unit $a$ modulo $15$. Exact modular-word closure
extends the known span from dimension $29$ to dimension $30$ in each
profile. All recorded linear relations, symmetry and isometry identities
are checked in $\mathbb Q(\zeta_{15})$.

For the rank-$38$ profile, the image of the individual basis vector
$e_6$ becomes known. Its restriction is $(0,0,0,0,1,0)$.
The known rows give
\[
 [A_1A_1Z_6:\mathbf1]
 =\sum_j\frac{S_{1j}^{\,2}S_{6j}}{S_{0j}}
 =\frac13,
\]
contradicting Verlinde integrality. The sum is recomputed independently
using rational polynomial arithmetic modulo
\[
 \Phi_{15}(x)=x^8-x^7+x^5-x^4+x^3-x+1.
\]
Conjugating preserves the rational value $1/3$, so both twist branches
are excluded. Only the rank-$35$ profile remains.

For rank $35$, let $W$ be the resulting known $30$-dimensional modular
subspace. Its orthogonal complement is explicitly computed and verified
to be $T$-stable. Its five distinct $T$-eigenvalues are
\[
 \{\zeta_{15}^{2},\zeta_{15}^{3},\zeta_{15}^{8},
   \zeta_{15}^{10},\zeta_{15}^{12}\}.
\]
The center's anomaly-one representation factors through
$SL_2(\mathbb Z/15\mathbb Z)$ by the congruence theorem
\cite[Theorem~II and Remark~4.5]{DLN}. The Chinese remainder theorem
identifies this finite group with
$SL_2(\mathbb F_3)\times SL_2(\mathbb F_5)$.
The computation constructs all $24$ and $120$ matrices in these groups,
their conjugacy classes and the complete class-algebra multiplication
tensors. A class-algebra element with simple spectrum yields every
primitive central character; character orthogonality gives its degree.
Discrete Fourier inversion on powers of the upper unipotent gives its
exact $T$-spectrum. The computation verifies simple spectrum and all
integer degrees and eigenmultiplicities, so no irreducible is omitted.

Tensoring these complete irreducible lists and comparing spectra leaves
exactly two decompositions of $W^\perp$. Their dimensions are
$1+4$ and $1+2+2$. For the first decomposition $-I$ acts trivially; for
the second its negative eigenspace has dimension four. Charge conjugation
can interchange only simples with the same restriction column and twist.
The retained profile has exactly two such pairs, namely
$\{21,22\}$ and $\{23,24\}$, and all other classes are singletons.
An involution on these classes has negative eigenspace dimension at most
two. Thus the $1+2+2$ representation is impossible.

For completeness, the remaining irreducible block is constructed
directly, rather than inferred from approximate character values.
Put $z_p=\zeta_{15}^{15/p}$ and
$g_p=\sum_{x=0}^{p-1}z_p^{x^2}$. Two exact factors are
\[
 S_3=\frac1{g_3}\begin{pmatrix}1&2\\1&-1\end{pmatrix},
 \qquad T_3=\operatorname{diag}(1,z_3),
\]
and
\[
 (S_5)_{xy}=\frac{z_5^{-2xy}-z_5^{2xy}}{g_5}
 \quad(x,y\in\{1,2\}),\qquad
 T_5=\operatorname{diag}(z_5,z_5^4).
\]
Enumerating every Cayley-graph edge checks that these matrices define
representations of the stated finite groups. Their distinct $T$
eigenvalues and nonzero off-diagonal $S$ entries establish irreducibility.
The four-dimensional tensor product has the required spectrum
$\{2,3,8,12\}$ in exponent notation. The remaining one-dimensional
character has $S=1$ and $T=\zeta_{15}^{10}$.

Each complementary $T$-eigenspace is one-dimensional. Hence any other
realization of the surviving representation with this fixed $T$ differs
by a diagonal intertwiner. Symmetry of $S$ fixes the relative squares
of the four diagonal scalars, leaving precisely eight choices of three
signs after removing a common scalar. The reconstruction checks
$S=S^{\mathsf T}$, $S^2=I$ and $STS=T^{-1}ST^{-1}$ for every choice.
Four choices fail nonnegative integral Verlinde multiplicities. The
other four pass all $\binom{36}{3}=7140$ unordered nonunit triples,
and unit triples follow from the unit row and $S^2=I$.
Exact simultaneous row--column permutations show that the four surviving
matrices differ only by the two pair interchanges above. Complex
conjugation supplies the other twist branch. This proves the stated
exhaustive reduction.
\end{proof}

The complete $35\times35$ matrices and $35^3$ fusion tensor are retained
in \filename{research/round5/m4\_629/rank35\_modular\_completions.json}
and \filename{rank35\_fusion\_data.json}. Every nonunit basis object
generates this entire necessary fusion ring. The supplied restriction
matrix is a based-ring homomorphism, verified on all $3780$ unordered
product coordinates. These properties do not supply an associator.

The exact reduction certificate is
\filename{research/round5/m4\_629/modular\_reduction\_certificate.json}.
The script \filename{replay\_modular\_chain.py} reconstructs all three
bounded searches, the twist lists, modular closure, both Verlinde
calculations, the complete finite-group character calculation and every
modular completion. It runs with numerical optimizers disabled and is
also replayed under Python's \texttt{-O} option. The replay records
source and artifact hashes. Numerical infeasibility reports and earlier
search interruptions play no part in the proof.


\subsection{The multiplicity-five frontier}
\label{sec:frontier-1204}

\begin{proposition}\label{prop:1204-finite-profiles}
Every complex categorification of $\RR{5}{1204}$ is spherical and has a
pseudo-unitary Galois conjugate. Its central induction matrix belongs,
up to permutation of central simples, to an explicitly certified list of
$171$ restriction profiles. Their central ranks lie between $22$ and $51$.
The list is a necessary condition; it supplies neither a categorification
nor a complete fusion ring or modular datum for the center.
\end{proposition}

\begin{proof}
The unrestricted signed-sphericalization test of Section~\ref{sec:signed}
enumerates $512$ associative signed quotient tables, partitioned into
$16$ sign-gauge orbits. One orbit consists of gauges of the original
fusion tensor. The exact dimension-character tests reject all $46$
character components on the other orbit representatives. As in the
sphericality argument of Proposition~\ref{prop:frontier-galois-pu}, a
categorification must therefore admit a spherical structure. Among the
character components of the original ring, the only one without a zero
simple dimension is the Frobenius--Perron orbit. Conjugating to its
positive embedding gives pseudo-unitarity. The positive global dimension
is $336+192\sqrt3$.

We may consequently apply the central-induction criteria with positive
dimensions. Exhaustive low-column enumeration gives $39$ frames. Only
two admit the required character assignments. Both assignments on the
first frame are excluded by exact full-character twist equations. The
second frame has three identical low restriction columns; its six
assignments are precisely the permutations of those columns. Relabelling
the corresponding low simples identifies these six necessary-data
problems, so one representative suffices.

For this representative, write the restriction Gram equations, joint
Galois multiplicity equations, and low cubic-word equations as an
integer nonnegative system $An=b$. Here $n$ includes both high-column
counts and low fusion multiplicities. The supplement records $165$
exact bound certificates. Each upper bound $n_j\le B$ is proved by a
separating functional for the augmented system
\[
 An=b,\qquad n_j-s=B+1,\qquad n,s\ge0;
\]
its infeasibility and integrality exclude $n_j\ge B+1$. Replacing $n_j$
by $-n_j$ gives lower bounds in the same way. Every bound is checked
against the preceding certified system. Bounded affine-lattice
enumeration then visits $293972$ nodes and produces exactly $39800$
high restriction profiles. Requiring integral nonnegative low-word
multiplicities through degree six and nonnegative restriction remainders
leaves $1833$ profiles.

For all unordered pairs of nonunit low simples, apply
Theorem~\ref{thm:joint-restriction-word-gram}. The profiles determine
$43$ distinct joint restriction--word Gram matrices. The exact fixed-column
Schur bounds of Lemma~\ref{lem:exact-rank-one-gram} enumerate every possible
joint coefficient column. The resulting count, Gram, and restriction
systems exclude $1568$ profiles by exact separating functionals. The
height-three full-character twist systems exclude another $91$.
Two additional joint count systems have inconsistent parity equations.
Finally, adjoining all cubic words in the last three low simples gives
an exact joint Gram obstruction for one further profile. Thus precisely
\[
 1833-1568-91-2-1=171
\]
profiles remain in the certified necessary domain. None is discarded
because a numerical optimizer reports infeasibility or exhausts its
running-time budget.

The complete certificate is
\filename{research/round5/m5_1204/profile_certificate.json}; the sparse
retained matrices are in
\filename{research/round5/m5_1204/remaining_profiles.json}.
The replay program \filename{software/round5_1204.py}, with its exact
bounded integer enumerator \filename{software/round5_1204.cpp}, reconstructs
all low frames, the labelling quotient, every integral bound, the complete
profile enumeration, and every joint-column domain before checking the
separating and congruence certificates. Numerical optimization is disabled
during replay, including under Python's optimized execution mode.
\end{proof}


\subsection{The remaining multiplicity-eight ring}
\label{subsec:frontier-1660}

The ring $R_{6,8,1660}$ is self-dual and contains the Ising based
subcategory with $g=b_1$ and $\sigma=b_4$. The unrestricted signed
reduction recalled above reduces every categorification to the
pseudo-unitary case. Its complete replay is
\filename{research/round5/m8_1660_conductor/replay_reduction.py}.
We use positive dimensions and put $D=\operatorname{FPdim}(\mathcal C)$.

\begin{proposition}[Actual-twist support and individual capacity]
\label{prop:1660-actual-support}
For a central simple with restriction $v=(v_0,\ldots,v_5)$, put
\[
 A(v)=v_0+v_1+v_2+v_3+2v_4-2v_5,\qquad
 B(v)=v_0+v_1-v_2-v_3+v_4+v_5.
\]
Then $A(v)\ne0$ implies $\theta_v^{24}=1$, and $B(v)\ne0$
implies $\theta_v^{12}=1$. Every actual central twist has order
dividing $48$.

Every high column containing $\sigma$ satisfies
\[
 v_0=v_1=0,\qquad v_4=1,\qquad v_2+v_3=4,\qquad v_5=3.
\]
There are exactly eight such central simples, each of dimension
$D\sqrt2/8$. Each of the four allowed actual Ising twists occurs on
exactly two of these simples.
\end{proposition}
\begin{proof}
The two saturated dimension-field characters have conductors $24$ and
$12$. Evaluating their Fourier coefficients in
\eqref{eq:saturated-column-support} gives nonzero constant multiples
of $A$ and $B$, respectively. The character representatives are
$1,5,7,11$; their sign vectors are $(1,1,-1,-1)$ and
$(1,-1,-1,1)$. The exact coefficient calculation is reconstructed by
\filename{software/dimension_block_support.py}.
Proposition~\ref{prop:saturated-conductor-support} gives the order bounds.

If a high column contains $\sigma$, its actual twist is primitive of
order $16$ and its $g$-coordinate is zero, by
Corollary~\ref{cor:fibonacci-ising-projection}. Therefore both Fourier
coefficients vanish. Solving these two integer linear equations gives
$v_2+v_3=4v_4$ and $v_5=3v_4$. Exact evaluation of the dimensions gives
\[
 \frac{v_4\operatorname{FPdim}(v)}{D\sqrt2/4}
       =\frac{v_4^2}{2}.
\]
The individual capacity bound \eqref{eq:individual-projection-capacity}
makes this ratio at most one. Hence $v_4=1$.
The induction Gram entry at $(\sigma,\sigma)$ is eight, and no low
column contains $\sigma$. Thus there are eight such simples, each
contributing half of one root mass, proving the stated allocation.
This uses individual actual-root masses before relative trace projection.

For the exponent bound, consider a high column on which both $A$ and
$B$ vanish. If $v_4>0$, the Ising support gives order $16$.
Otherwise $v_5=2(v_0+v_1)$ and $v_2+v_3=3(v_0+v_1)$.
Since $v_0=0$ and the column is nonzero, $v_1>0$, so its twist is $-1$.
All other high columns have twist order dividing $24$, and low twists
are one. Their least common multiple divides $48$.
\end{proof}

\begin{proposition}[Complete necessary restriction profiles]
\label{prop:1660-profiles}
For a categorification of $R_{6,8,1660}$, the pseudo-unitary Galois
conjugate has one of two low frames, each with two admissible character
labellings. After the actual-twist support and capacity bounds,
the complete induction Gram equations, joint-character and
normalized-dimension Galois balance, integer cubic Verlinde
multiplicities, and nonnegative low-product restriction remainders,
the numbers of high restriction profiles are as follows, per labelling:
\[
\begin{array}{c|r|r|c}
 \text{low frame}&\text{cubic constraints}&\text{words through degree six}
 &\text{range of possible center ranks}\\\hline
 0&371&371&25\text{--}76\\
 1&2011&1999&25\text{--}82
\end{array}
\]
For a fixed low frame the two labellings give the same surviving
restriction profiles. Every retained profile also passes entrywise
nonnegativity of the residual Gram matrix of all fifteen quadratic
words in the five nonunit low simples.
\end{proposition}
\begin{proof}
Enumerate the high columns using the induction Gram bounds,
formal-codegree divisibility, complete character integrality and
Galois closure, and the Ising projection constraints. Apply
Proposition~\ref{prop:1660-actual-support} and
Lemma~\ref{lem:simple-restriction-support}; the latter bounds the pure
$g$ multiplicity by one. There remain $20$ high columns in the first
frame and $52$ in the second.

Their nonnegative integral multiplicities satisfy a bounded affine
system. The Gram and Galois equations have ranks $11$ and $30$;
exact row reduction leaves $9$ and $22$ free coordinates. Diagonal
Gram entries bound every multiplicity. Cubic character formulas give
integrality congruences and nonnegativity inequalities. Restricting
each low product to the source based ring supplies coordinate bounds.

The bounded free coordinates are enumerated with exact interval and
divisibility pruning. The complete searches visit $1009$ and $63366$
nodes, respectively, per character labelling. Higher-word tests use
\eqref{eq:unified-word-pairing}; degrees four through six remove
twelve profiles in the second frame and none in the first.
The quadratic-word residual Gram test subtracts the known low
constituents as in Theorem~\ref{thm:joint-restriction-word-gram}.

All domains, rational row reductions, pruning decisions and retained
profiles are reconstructed by \filename{software/round5_1660.py}.
The records are \filename{research/round5/m8_1660/low0_profiles.json}
and \filename{low1_profiles.json} in the same directory.
The source-bound receipts \filename{low0_replay.json} and
\filename{low1_replay.json} record complete optimizer-free replays
under \texttt{python -O}. Consequently neither numerical solver
status nor Python assertions are used to prove completeness.
\end{proof}

The eight $\sigma$-supported simples have only two possible restriction
multisets in these lists. In frame zero all eight have column
$(0,0,2,2,1,3)$. In frame one, $1628$ profiles have this same multiset;
the other $371$ have six copies of this column and one each of
$(0,0,1,3,1,3)$ and $(0,0,3,1,1,3)$, per labelling.
The exact first-trace cancellation condition on opposite actual Ising
roots admits an allocation for both multisets, so it supplies no
additional exclusion by itself. The exact control
\filename{research/round5/m8_1660/verify_sigma_partitions.py}
reconstructs these multisets and their allocations, and checks the
displayed Fourier forms on all six basis columns.

\begin{proposition}[Complete actual-root profile reduction]
\label{prop:1660-actual-frontier}
The actual-root equations leave precisely the following necessary
profile alternatives:
\[
\begin{array}{c|c|r|r}
\text{low frame}&\text{character labelling}&\text{input profiles}
 &\text{remaining profiles}\\\hline
0&(0,5,4,1,3,2)&371&168\\
1&(0,5,4,1,2,3)&1999&798\\
1&(0,5,4,1,3,2)&1999&798
\end{array}
\]
All three branches have second indicators $(1,1,1,-1,1)$.
Every retained profile has eight copies of the balanced
$\sigma$-supported column $(0,0,2,2,1,3)$.
Their possible center ranks are $28,31,\ldots,76$ in the first frame
and $28,31,\ldots,82$ in the second.
\end{proposition}
\begin{proof}
The proved exponent bound permits the use of the $48$ actual roots as
the complete twist domain. Impose the induction Gram equations,
first and second indicator traces, cyclic-rotation multiplicities,
joint-character reciprocity, normalized-dimension Galois balance,
the four individual Ising-root masses, and the low-product
restriction constraints. There are sixteen low-frame, character,
and indicator branches. Exact Farkas vectors exclude thirteen of
them. In the first retained branch, three shared separating vectors
exclude $203$ of the $371$ restriction profiles.

For each second-frame branch the shared separating vectors leave
$966$ profiles. A six-row integer combination of the fixed-profile
system has every coefficient divisible by four, but its right-hand
side is not divisible by four for any asymmetric $(1,6,1)$
$\sigma$-pattern. This excludes all such patterns, including the
$168$ that remained after the separating-vector tests, leaving
$798$ profiles per labelling. The checker intersects these results
with the complete degree-six survivor indices; it does not truncate
the stored $2011$-element cubic list.

The compact record
\filename{research/round5/m8_actual48/certificate.json} stores the
shared integer vectors, congruences, and complete remaining index
lists. Its replay program \filename{certificate.py} reconstructs
all sixteen systems and checks every applicable profile exactly.
The source-bound \filename{replay_O.json} records a complete run
with numerical optimizers and Python assertions disabled.
\end{proof}

\begin{proposition}[A forced Ising-center subrepresentation]
\label{prop:1660-ising-subrepresentation}
The canonical modular representation of a hypothetical center
contains the full nine-dimensional representation of the Ising
center as a subrepresentation. Its central projection matrix $B$
satisfies $BB^\dagger=2I_9$.
\end{proposition}
\begin{proof}
For every occupied column in the complete lists of
Proposition~\ref{prop:1660-profiles}, its Ising projection is one of
\[
 (0,0,0),\quad(1,0,0),\quad(1,1,0),\quad(0,1,0),\quad(0,0,1).
\]
At the allowed actual twist, each nonzero projection is a single
simple of the Ising center, with multiplicity one. This finite
support assertion is reconstructed for all $4740$ labelled profiles
by \filename{research/round5/m8_modrep/ising_complement.py}.
Corollary~\ref{cor:disjoint-central-projection} therefore applies:
the rows of $B$ have disjoint supports, their individual positive
dimension masses make them nonzero, and $BB^\dagger$ is a positive
scalar matrix. The four $\sigma$ rows each contain exactly two
objects, so that scalar is two.
\end{proof}

The finite-quotient criterion of
Subsection~\ref{subsec:modular-word-closure} is implemented for
$\mathrm{SL}_2(\mathbb Z/48\mathbb Z)$ using its two Chinese-remainder
factors. The exact class algebras give $76$ irreducible characters
for $\mathrm{SL}_2(\mathbb Z/16\mathbb Z)$ and seven for
$\mathrm{SL}_2(\mathbb Z/3\mathbb Z)$. Their $532$ tensor products
give $504$ distinct signed $T$-spectra. The independent replay checks
all class-algebra identities, character orthogonality, degrees,
central signs and exact Fourier multiplicities. Two Ising-center
spectra serve as positive controls. These records are in
\filename{research/round5/m8_modrep/}.

One may apply this spectrum cone to the complement of the forced
Ising-center representation. Its known $T$-spectrum has multiplicity
three at $1$, two at $-1$, and one at each of the four Ising roots.
The signed-trace lower bounds must also subtract these multiplicities:
they become $3$ at $1$, $-2$ at $-1$, and $-1$ at the four Ising roots.
The complement itself need not have permutation duality.

The full-spectrum and forced-complement searches produced no further
exact exclusion. The finite-group character enumeration is complete;
the coupled nonnegative integral completion problem for the remaining
profiles is not solved. The ring remains unresolved. The remaining
lists are necessary alternatives, not induction matrices or modular
data asserted to come from fusion categories.

